\section{Definition of Correlation functions and free energies}

In all this section, the curve $\curve(x,y)=0$ and a symmetric matrix $\kappa$ are given and fixed.
The unfamiliar reader may choose $\kappa=0$ since most usual applications (matrix models) correspond to that case.


\subsection{Notations}


Consider an arbitrary point $p\in\bar\Sigma$, and a point $q$ of $\bar\Sigma$ which is in the vicinity of a branch point $a_i$ (so that $\qbar$ is well defined).
We define:
\bd\label{defdiagrule}
 Diagrammatic rules:
\beq
{\mathbf{\rm vertex:}}\qquad \om(q) = (y(q)-y(\qbar))dx(q)
\eeq
\beq
{\mathbf{\rm line-propagator:}}\qquad B(p,q)
\eeq
\beq
{\mathbf{\rm arrow-propagator:}}\qquad dE_{q}(p) = {1\over 2} \int_{q}^{\qbar} B(\xi,p)
\eeq
where the integration path is a  path which lies entirely in a vincinity of $a_i$ (thus it is uniquely defined)\footnote{This definition is the opposite
of the notation used in \cite{eyno,CEO} since the integral goes from $q$ to $\qbar$ instead of going from
$\qbar$ to $q$.}.
\beq
{\mathbf{\rm root:}}\qquad \Phi(q) = \int_{o}^q y dx
\eeq
where $o$ is an arbitrary base point on the curve, i.e. $\Phi$ is an arbitrary antiderivative of $y dx$, i.e.  $d\Phi=ydx$.
\ed
The reason why we call these objects diagramatic rules and vertices, propagator or root, is explained in section \ref{sectiondiagrepresent} below.

Notice that $dE$ depends on $i$, i.e. on which branchpoint we are considering, but we omit to mention the index $i$ in order to make the notations easier to read.
In all what follows it is always clear which $i$ is being considered.

\bigskip
{\bf Notation for subset of indices:}

Given a set of points of the curve $\{ p_{1}, p_{2},\dots, p_{n}\}$, 
if $K=\{ i_1,i_2,\dots,i_k\}$ is any subset of $\{ 1,2,\dots, n\}$, we denote:
\beq\label{notationsubsetint}
{\mathbf p_K} = \{ p_{i_1}, p_{i_2},\dots, p_{i_k}\}
\eeq


\subsection{Correlation functions and free energies}

\subsubsection{Correlation functions}

The {\bf $k$-point correlation functions to order $g$, $W_{k}^{(g)}$}, 
are meromorphic multilinear forms, defined
by the following recursive triangular system:
\bd\label{defloopfctions}
Correlation functions
\beq
W_k^{(g)}=0 \quad {\rm if}\,\, g<0
\eeq
\beq
W_1^{(0)}(p) = 0
\eeq
\beq\encadremath{
W_2^{(0)}(p_1,p_2) = B(p_1,p_2)
}\eeq
and define recursively (remember that ${\bf p_K}$ is a $k$-uplet of points cf eq.\ref{notationsubsetint}):
\beq\label{defWkgrecursive}\encadremath{
\begin{array}{l} 
W_{k+1}^{(g)}(p,{\bf p_K}) 
= \Res_{q\to {\bf a}} {dE_{q}(p)\over \om(q)}\,\Big(  \cr
\qquad \sum_{m=0}^g \sum_{J\subset K} W_{|J|+1}^{(m)}(q,{\bf p_J})W_{k-|J|+1}^{(g-m)}(\qbar,{\bf p_{K/J}})
+ W_{k+2}^{(g-1)}(q,\qbar,{\bf p_K}) \Big) \cr
\end{array}
}\eeq
\ed
This system is triangular because all terms in the RHS have lower $2g+k$ than in the LHS and
given $W_1^{(0)}$ and $W_2^{(0)}$, it has a unique solution.


Notice that $W_{k+1}^{(g)}(p,p_1,\dots,p_k)$ is a multilinear meromorphic form in each of its arguments, it is clearly symmetric in the last $k$-ones, and we prove below (theorem \ref{thsymWk}) that it is in fact symmetric in all its arguments.

More properties of $W_{k+1}^{(g)}$ are studied below in section \ref{propcorrel}.

\subsubsection{Free energies}

We define the {\bf free energies} which are complex numbers:

\bd\label{defFg}
Free energies.

For $g>1$
\beq\encadremath{
F^{(g)} = {1\over 2-2g}\,\sum_i \Res_{q\to a_i} \Phi(q) W_{1}^{(g)}(q)
}\eeq
and
\beq\encadremath{
F^{(1)}=-{1\over 2}\ln{(\tau_{Bx})}\,\,-{1\over 24}\ln{\left(\prod_i y'(a_i) \right)}- \ln\left(\det \kappa\right)
}\eeq
where
\beq
y'(a_i) = {dy(a_i)\over dz_i(a_i)}
\virg
z_i(p) = \sqrt{x(p)-x(a_i)}
\eeq
and $\tau_{Bx}$ is the Bergmann $\tau$-function defined in \eq{deftauBx}.

$F^{(0)}$ is defined in the next section.

\ed

\subsubsection{Leading order free energy  \texorpdfstring{$F^{(0)}$}{F0}.}

Let us define $F^{(0)}$ as follows
:
%by one of these two quantities

\beq\label{defF0}\encadremath{
F^{(0)}
= {1\over 2}\sum_{\alpha} \Res_\alpha V_\alpha ydx + {1\over 2} \sum_{\alpha} t_\alpha \mu_\alpha
- {1\over 4i\pi}\sum_i \oint_{\acycle_i} ydx \oint_{\bcycle_i} ydx 
%&= &  -{1\over 2}\sum_{\alpha,\beta} \Res_{p\to \alpha}\Res_{q\to \beta} V_\alpha(p) B(p,q) V_\beta(q)
%+ \sum_{\alpha,\beta} t_\beta \Res_\alpha V_\alpha dS_{\beta,o} \cr
%&& - {1\over 2}\sum_{\alpha, \beta} t_\alpha t_\beta \ln{(\gamma_{\alpha,\beta})}
%- {1\over 2}\epsilon^t \Gamma  \cr
}\eeq
where $\mu_\alpha$ is given by 
%any of the equivalent following expressions:
\beq
\mu_\alpha
%&=& \sum_\beta \Res_\beta V_\beta dS_{\alpha,o}
%- \sum_{\beta} t_\beta  \ln{(\gamma_{\alpha,\beta})} - 2i\pi \sum_i \epsilon_i (u_i(\alpha)-u_i(o)) \cr
= \int_\alpha^{o} (ydx - dV_\alpha + t_\alpha {dz_\alpha\over z_\alpha}) + V_\alpha(o) - t_\alpha \ln{(z_\alpha(o))} 
\eeq
%and where
%\beq
%\ln{\gamma_{\alpha,\alpha}} = -\int_\alpha^o (dS_{\alpha,o'} + {dz_\alpha\over z_\alpha}) +  \ln{(z_\alpha(o))}
%\eeq
%\beq
%\ln{\gamma_{\alpha,\beta}} = \ln{\left(\primef(\alpha,\beta)\primef(o,o')\over \primef(\alpha,o')\primef(\beta,o)\right)}
%\eeq
%as well as
%\beq
%\Gamma  := \oint_{\bcycle} ydx.
%\eeq
Notice that $\mu_\alpha$ depends on some base point $o$, but the sum $\sum_\alpha t_\alpha \mu_\alpha$ does not.

\subsubsection{Special free energies and correlation functions}
\label{sectdefspfree}

All the quantities defined so far, were defined with the $\kappa$-modified cycles and modified Bergmann kernel.
Let us also define them for $\kappa=0$ (for instance $F^{(1)}$ obviously needs another definition).

Therefore we also define the unmodified quantities corresponding to $\kappa=0$, as:
\beq
\label{defspcorr}
\forall k,g, \qquad 
\underline{W}_{k}^{(g)}(p_1,\dots,p_k)
:= \left. W_{k}^{(g)}(p_1,\dots,p_k)
\right|_{\kappa = 0},
\eeq
\beq \label{defspfree}
{\rm for}\,\, g>1, \qquad
\underline{F}^{(g)} := {1\over 2-2g}\,\sum_i \Res_{q\to a_i} \Phi(q) \underline{W}_{1}^{(g)}(q)
= \left.
F^{(g)} \right|_{\kappa = 0} \,\, ,
\eeq
\beq
\underline{F}^{(1)}=-{1\over 2}\ln{(\tau_{\underline{B}x})}\,\,-{1\over 24}\ln{\left(\prod_i y'(a_i) \right)} \,\, ,
\eeq
\beq
{\rm and}\qquad
\underline{F}^{(0)}
= {1\over 2}\sum_{\alpha} \Res_\alpha V_\alpha ydx + {1\over 2} \sum_{\alpha} t_\alpha \mu_\alpha
- {1\over 4i\pi}\sum_i \oint_{\underline\acycle_i} ydx \oint_{\underline\bcycle_i} ydx \, .
\eeq

\br
The special functions, except $\underline{F}^{(1)}$ and $\underline{F}^{(0)}$, are obtained by changing $B$ and $dS$ by
$\Bergmann$ and $\underline{dS}$ in the diagrammatic rules defined in section \ref{sectiondiagrepresent}.
\er

\subsection{Tau function}

\bd
We define the {\bf tau-function} as the formal power series in $N^{-2}$:
\beq\encadremath{
\ln{(Z_N(\curve))} = -\sum_{g=0}^\infty N^{2-2g} F^{(g)}.
}\eeq
\ed
We show in section \ref{sectintegrability} that $Z_N(\curve)$ is indeed a tau-function because it obeys Hirota equations, order by order in $N^{-2}$.


\subsection{Properties of correlation functions}\label{propcorrel}


The loop functions defined in definition \ref{defloopfctions} satisfy the following theorems, whose {\bf proofs can be found
in Appendix \ref{appprop}}:


%---------------------------------thpolesW30------------------
\bt\label{thW30}
The correlation function $W_3^{(0)}$ is worth:
\beq
W_3^{(0)}(p,p_1,p_2) = \Res_{q\to \bfa} {B(q,p)B(q,p_1)B(q,p_2)\over dx(q) dy(q)}
\eeq
\et
In particular, $W_3^{(0)}$ is symmetric in its 3  variables.


%---------------------------------thpolesWkgbp------------------
\bt\label{thpolesWkgbp}
For $(k,g)\neq (1,0)$, the loop function $W_{k+1}^{(g)}(p,p_1,\dots,p_k)$ has poles (in any of its variables $p,p_1,\dots,p_k$) only at the branch points.
%, of degree $6g+2k-4$.
\et


%---------------------------------thWkcycle------------------
\bt\label{thWkcycle}
For every ${\cal A}$ cycle we have:
\beq
\forall i=1,\dots,\genus \qquad \oint_{p\in{\cal A}_i} W_{k+1}^{(g)}(p,p_1,\dots,p_k) = 0
,
\eeq
\beq
\forall i=1,\dots,\genus, \forall m=1,\dots,k \qquad \oint_{p_m\in{\cal A}_i} W_{k+1}^{(g)}(p,p_1,\dots,p_k) = 0
.
\eeq
\et


%---------------------------------thsumWk------------------
\bt\label{thsumWk}
For every $k$ and $g$, we have:
\beq\label{thsumWk1}
\sum_i {W_{k+1}^{(g)}(p^i,p_1,\dots,p_k) \over dx(p^i)}= \delta_{k,1}\delta_{g,0}\,{dx(p_1)\over (x(p)-x(p_1))^2}
\eeq
and if $k\geq 1$:
\beq\label{thsumWk2}
\sum_i {W_{k+1}^{(g)}(p_1,p^i,p_2,\dots,p_k) \over dx(p^i)} = \delta_{k,1}\delta_{g,0}\,{dx(p_1)\over (x(p)-x(p_1))^2}
\eeq
where we recall that $p^i$ are all the points such that $x(p^i)=x(p)$ (see section \ref{sectgeomalg}).
\et


%---------------------------------thPkpol------------------
\bt\label{thPkpol}
For $(k,g)\neq (0,1)$,
\bea
P_k^{(g)}(x(p),{\bf p_K})
&=& {1\over dx(p)^2}\,\sum_i \Big[-2y(p^i)dx(p) W_{k+1}^{(g)}(p^i,{\bf p_K}) + W_{k+2}^{(g-1)}(p^i,p^i,{\bf p_K}) \cr
&& \qquad \qquad + \sum_{m=0}^g \sum_{J\subset K} W_{j+1}^{(m)}(p^i,{\bf p_J})W_{k-j+1}^{(g-m)}(p^i,{\bf p_{K/J}}) \Big] \cr
\eea
 is a rational function of $x(p)$, with no poles at branch-points.
\et



%---------------------------------thsymWk------------------
\bt\label{thsymWk}
$W_k^{(g)}$ is a symmetric function of its $k$ variables.
\et



\bc\label{corResxyWk}
\beq
\forall i,\qquad \Res_{p\to a_i} W^{(g)}_{k+1}(p,p_1,\dots,p_k) =0
,
\eeq
\beq
\forall i,\qquad \Res_{p\to a_i} x(p) W^{(g)}_{k+1}(p,p_1,\dots,p_k) =0
,
\eeq
\beq
\sum_i \Res_{p\to a_i} y(p) W^{(g)}_{k+1}(p,p_1,\dots,p_k) =0
,
\eeq
\beq
\sum_i \Res_{p\to a_i} x(p)y(p) W^{(g)}_{k+1}(p,p_1,\dots,p_k) =0
.
\eeq
\ec






%---------------------------------thintPhi------------------
\bt\label{thintPhi}
For $k\geq 1$ we have:
\beq
\Res_{p_{k+1}\to \bfa,p_1,\dots,p_k} \Phi(p_{k+1}) W^{(g)}_{k+1}({\bf p_K},p_{k+1}) = (2g+k-2) W^{(g)}_k({\bf p_K}) + \delta_{g,0}\delta_{k,1} y(p_1)dx(p_1)
.
\eeq
\et
Notice that for $k=0$ and $g\geq 2$, it holds by definition if we define $W^{(g)}_0= - F^{(g)}$.



\subsection{Diagrammatic representation \label{sectiondiagrepresent}}



The recursive  definitions of $W_k^{(g)}$ and $F^{(g)}$ can be represented {\bf graphically}.

We represent the multilinear form $W_k^{(g)}(p_1,\dots,p_k)$ as a blob-like ``surface'' with $g$ holes and
$k$ legs (or punctures) labeled with the variables $p_1,\dots, p_k$, and $F^{(g)}$ with $0$ legs and
$g$ holes.
\beq
W_{k+1}^{(g)}(p,p_1,\dots,p_k):=\begin{array}{r}
{\epsfxsize 4.5cm\epsffile{Wk.eps}}
\end{array}
\virg
F^{(g)}:= \begin{array}{r}
{\epsfxsize 2.5cm\epsffile{Fg.eps}}
\end{array}
\eeq

We represent the Bergmann kernel $B(p,q)$ (which is also $W_2^{(0)}$, i.e. a blob with 2 legs and no hole) as a straight non-oriented line between $p$ and $q$
\beq
B(p,q):= \begin{array}{r}
{\epsfxsize 2cm\epsffile{berg.eps}}
\end{array}.
\eeq

We represent ${dE_{q}(p)\over \om(q)}$ as a straight arrowed line with the arrow from $p$ towards $q$, and with a tri-valent vertex whose legs are $q$ and $\qbar$
\beq
{dE_{q}(p)\over \om(q)}:= \begin{array}{r}
{\epsfxsize 3cm\epsffile{vertex.eps}}
\end{array}.
\eeq

\medskip

\subsubsection*{Graphs}

\bd\label{defgraphs}
For any $k\geq 0$ and $g\geq 0$ such that $k+2g\geq 3$, we define:

Let ${\cal G}_{k+1}^{(g)}(p,p_1,\dots,p_k)$ be the set of connected trivalent graphs defined as follows:
\begin{enumerate}

\item there are $2g+k-1$ tri-valent vertices called vertices.
\item there is one 1-valent vertex labelled by $p$, called the root.
\item there are $k$ 1-valent vertices labelled with $p_1,\dots,p_k$ called the leaves.
\item There are $3g+2k-1$ edges.
\item Edges can be arrowed or non-arrowed. There are $k+g$ non-arrowed edges and $2g+k-1$ arrowed edges.
\item The edge starting at $p$ has an arrow leaving from the root $p$.
\item The $k$ edges ending at the leaves $p_1,\dots, p_k$ are non-arrowed.
\item The arrowed edges form a "spanning\footnote{It goes through all vertices.} planar\footnote{planar tree means that the left child and right child are not equivalent. The right child is marked by a black disk on the outgoing edge.} binary skeleton\footnote{a binary skeleton tree is a binary tree from which we have removed the leaves, i.e. a tree with vertices of valence 1, 2 or 3.} tree" with root $p$. The arrows are oriented from root towards leaves. In particular, this induces a partial ordering of all vertices.
\item There are $k$ non-arrowed edges going from a vertex to a leaf, and $g$ non arrowed edges joining two inner vertices. Two inner vertices can be connected by a non arrowed edge only if one is the parent of the other along the tree.
\item If an arrowed edge and a non-arrowed inner edge come out of a vertex, then the arrowed edge is the left child. This rule
only applies when the non-arrowed edge links this vertex to one of its descendants (not one of its parents).

\end{enumerate}

\ed


We have the following useful lemma:
\bl\label{lemgraphaddleg}
There is a $1$ to $3g+2k-1$ map from ${\cal G}_{k+1}^{(g)}(p,{\bf p_K})$ to ${\cal G}_{k+2}^{(g)}(p,{\bf p_K},p_{k+1})$.
\el
\proof{If $G$ is  a graph in ${\cal G}_{k+2}^{(g)}(p,p_1,\dots,p_k,p_{k+1})$, remove the non-arrowed edge attached to the leaf $p_{k+1}$ and remove the corresponding vertex, and merge the incoming and the other outgoing edges of that vertex. You clearly get a graph $G'\in{\cal G}_{k+1}^{(g)}(p,p_1,\dots,p_k)$. It is clear that the same graph is obtained $3g+2k-1$ times (the number of edges of $G'$.
And it is clear that from any $G'\in{\cal G}_{k+1}^{(g)}(p,p_1,\dots,p_k)$, you can obtain $3g+2k-1$ graphs $G\in{\cal G}_{k+2}^{(g)}(p,p_1,\dots,p_k,p_{k+1})$ by adding a new vertex on any edge, and linking this new vertex to the leaf $p_{k+1}$.}

\vs

\noindent {\bf Example of ${\cal G}_{1}^{(2)}(p)$}

As an example, let us build step by step all the graphs of ${\cal G}_{1}^{(2)}(p)$, i.e. $g=2$ and $k=0$.

We first draw all planar binary skeleton trees with one root $p$ and $2g+k-1=3$ arrowed edges:
\beq
\begin{array}{r}
{\epsfxsize 2cm\epsffile{w12.eps}}
\end{array}
\virg
\begin{array}{r}
{\epsfxsize 2.5cm\epsffile{w21.eps}}
\end{array}.
\eeq
%
Then, we draw $g+k=2$ non-arrowed edges in all possible ways such that every vertex is trivalent, also satisfying rule 9) of definition.\ref{defgraphs}. There is only
one possibility for the first graph and two for the second one:
\beq
\begin{array}{r}
\begin{array}{r}
{\epsfxsize 2.5cm\epsffile{y123.eps}}
\end{array}
\virg
{\epsfxsize 2.2cm\epsffile{y122.eps}}
\end{array}
\virg
\begin{array}{r}
{\epsfxsize 3cm\epsffile{y121.eps}}
\end{array}.
\eeq

It just remains to specify the left and right children for each vertex. The only possibilities in accordance with rule 10) of def.\ref{defgraphs} are\footnote{ Note that the graphs are not
necessarily planar.}:
\bea
\begin{array}{r}
{\epsfxsize 2.2cm\epsffile{y123.eps}}
\end{array}
\virg
\begin{array}{r}
{\epsfxsize 2.5cm\epsffile{x124.eps}}
\end{array}
\virg
\begin{array}{r}
{\epsfxsize 3cm\epsffile{x122.eps}}
\end{array} \cr
\begin{array}{r}
{\epsfxsize 2.5cm\epsffile{x123.eps}}
\end{array}
\virg
\begin{array}{r}
{\epsfxsize 3cm\epsffile{x121.eps}}
\end{array} .\cr
\eea

In order to simplify the drawing, we can draw a black dot to specify the right child. This way one gets only planar graphs.
\bea
\begin{array}{r}
{\epsfxsize 2.2cm\epsffile{w122.eps}}
\end{array}
\virg
\begin{array}{r}
{\epsfxsize 2.5cm\epsffile{w125.eps}}
\end{array}
\virg
\begin{array}{r}
{\epsfxsize 3cm\epsffile{w123.eps}}
\end{array} \cr
\begin{array}{r}
{\epsfxsize 2.5cm\epsffile{w124.eps}}
\end{array}
\virg
\begin{array}{r}
{\epsfxsize 3cm\epsffile{w121.eps}}
\end{array} \cr
\eea
%
Remark that without the prescriptions 9) and 10), one would get 13 different graphs whereas we only have 5.


\subsubsection*{Weight of a graph}

Consider a graph $G\in {\cal G}_{k+1}^{(g)}(p,p_1,\dots,p_k)$.
Then, to each vertex $i=1,\dots,2g+k-1$ of $G$, we associate a label $q_i$, and we associate $q_i$ to the beginning of the left child edge, and $\qbar_i$ to the right child edge.
Thus, each edge (arrowed or not), links two labels which are points on the Riemann surface $\bar\Surf$.
\begin{itemize}
\item To an arrowed edge going from $q'$ towards $q$, we associate a factor ${dE_q(q')\over (y(q)-y(\qbar))dx(q)}$.
\item To a non arrowed edge going between $q'$ and $q$ we associate a factor $B(q',q)$.
\item Following the arrows backwards (i.e.  from leaves to root), for each vertex $q$, we take a residue at $q\to \bfa$, i.e. we sum over all branchpoints.
\end{itemize}
After taking all the residues, we get the weight of the graph:
\beq
w(G)
\eeq
which is a multilinear form in $p,p_1,\dots,p_k$.

Similarly, we define weights of linear combinations of graphs by:
\beq
w(\alpha G_1 + \beta G_2) = \alpha w(G_1) + \beta w(G_2)
\eeq
and for a disconnected graph, i.e. a product of two graphs:
\beq
w( G_1  G_2) = w(G_1)  w(G_2)
.
\eeq


\bt\label{thdiagrepr}
We have:
\beq
W_{k+1}^{(g)}(p,p_1,\dots,p_k) = \sum_{G\in {\cal G}_{k+1}^{(g)}(p,p_1,\dots,p_k)}\,\, w(G) = w\left(\sum_{G\in {\cal G}_{k+1}^{(g)}(p,p_1,\dots,p_k)}\,\, G\right)
\eeq
\et
\proof{This is precisely what the recursion equations \ref{defWkgrecursive} of def.\ref{defloopfctions} are doing.
Indeed, one can represent them diagrammatically by
\beq
\begin{array}{r}
{\epsfxsize 10cm\epsffile{recurW.eps}}
\end{array}.
\eeq
}

\bigskip


Such graphical notations are very convenient, and are a good support for intuition and even help proving some relationships.
It was immediately noticed after \cite{eynloop1mat} that those diagrams look very much like Feynman graphs, and there was a hope that they could be the Feynman's graphs for the Kodaira--Spencer theory.
But they ARE NOT Feynman graphs, because Feynman graphs can't have non-local restrictions like the fact that non oriented lines can join only a vertex and one of its descendent.

Those graphs are merely a notation for the recursive definition \ref{defloopfctions}.

\bl\label{lemmasymfactor} {\bf Symmetry factor:}
The weight of two graphs differing by the exchange of the right and left children of a vertex are the same. Indeed, the distinction between right and left child is just a way of encoding symmetry factors.
\el

\proof{
This property follows directly from theorem \ref{thsumWk} and the definition \eq{defWkgrecursive}.
Consider one term contributing to the first part of RHS of \eq{defWkgrecursive}:
\beq
\begin{array}{l}
\displaystyle \Res_{q \to {\bf a}} {dE_{q}(p) \over \omega(q)} W_{|J|+1}^{(m)}(q,{\bf P_J}) W_{k-|J|+1}^{(g-m)}(\qbar,{\bf P_{K/J}}) \cr
\displaystyle \qquad \quad =  - \Res_{q \to {\bf a}} {dE_{q}(p) \over \omega(q)} W_{|J|+1}^{(m)}(q,{\bf P_J}) W_{k-|J|+1}^{(g-m)}(q,{\bf P_{K/J}}) \cr
\displaystyle \qquad \quad =  \Res_{q \to {\bf a}} {dE_{q}(p) \over \omega(q)} W_{|J|+1}^{(m)}(\qbar,{\bf P_J}) W_{k-|J|+1}^{(g-m)}(q,{\bf P_{K/J}}) .\cr
\end{array}
\eeq
where the equalities are obtained by adding terms without residues at the branch points to the integrand and
using theorem \ref{thsumWk}. 
One can perform the same trick for the second term in \eq{defWkgrecursive} and this proves the lemma.}


\subsection{Examples.}

Let us present some examples of correlation functions and free energy for low order.

\subsubsection*{Correlation functions.}

To leading order, one has the first correlation functions given by:
\beq
W_{2}^{(0)}(p,q) = B(p,q).
\eeq

\bea
W_{3}^{(0)}(p,p_1,p_2)&=& \begin{array}{r}
{\epsfxsize 3.5cm\epsffile{w31.eps}}
\end{array}
+
\begin{array}{r}
{\epsfxsize 3.5cm\epsffile{w32.eps}}
\end{array} \cr
&=& \Res_{q \to {\bf a}} {dE_q(p) \over \omega(q)} \left[ B(q,p_1) B(\overline{q},p_2) + B(\overline{q},p_1) B(q,p_2) \right] \cr
\eea

\bea
W_4^{(0)}(p,p_1,p_2,p_3)
&=& 
\begin{array}{r}
{\epsfxsize 4cm\epsffile{W4.eps}}
\end{array}
+ \;\;\;\; \hbox{perm. (1,2,3)}\cr
&& + \begin{array}{r}
{\epsfxsize 4cm\epsffile{W41.eps}}
\end{array}
+ \;\;\;\; \hbox{perm. (1,2,3)}\cr
&=& \Res_{q \to {\bf a}} \Res_{r \to {\bf a}} {dE_q(p) \over \omega(q)} {dE_r(q) \over \omega(r)}
\left[ B(\qbar,p_1)B(r,p_2)B(\overline{r},p_3) \right. \cr
&& + B(\qbar,p_1)B(\overline{r},p_2)B(r,p_3) + B(\qbar,p_2)B(r,p_1)B(\overline{r},p_3) \cr
&& + B(\qbar,p_2)B(\overline{r},p_1)B(r,p_3) + B(\qbar,p_3)B(r,p_2)B(\overline{r},p_1) \cr
&& \left. + B(\qbar,p_3)B(\overline{r},p_2)B(r,p_1)\right] \cr
&& + \Res_{q \to {\bf a}} \Res_{r \to {\bf a}} {dE_{q}(p) \over \omega(q)} {dE_r(\qbar) \over \omega(r)}
\left[ B(q,p_1)B(r,p_2)B(\overline{r},p_3) \right. \cr
&& + B(q,p_1)B(\overline{r},p_2)B(r,p_3) + B(q,p_2)B(r,p_1)B(\overline{r},p_3)\cr
&& + B(q,p_2)B(\overline{r},p_1)B(r,p_3) + B(q,p_3)B(r,p_2)B(\overline{r},p_1) \cr
&& \left. + B(q,p_3)B(\overline{r},p_2)B(r,p_1)\right] \cr
\eea


First orders for the one point correlation function read:
\bea
W_{1}^{(1)}(p) &=& \begin{array}{r}
{\epsfxsize 4cm\epsffile{w11.eps}}
\end{array} \cr
&=& \Res_{q \to {\bf a}} {dE_q(p) \over \omega(q)} B(q, \qbar)\cr
\eea

\bea
W_1^{(2)}(p)&=&
\begin{array}{r}
{\epsfxsize 4cm\epsffile{w121.eps}}
\end{array}
+
\begin{array}{r}
{\epsfxsize 4cm\epsffile{w123.eps}}
\end{array} \cr
&& +
\begin{array}{r}
{\epsfxsize 4cm\epsffile{w124.eps}}
\end{array}
+
\begin{array}{r}
{\epsfxsize 4cm\epsffile{w125.eps}}
\end{array} \cr
&& + \begin{array}{r}
{\epsfxsize 3.1cm\epsffile{w122.eps}}
\end{array}
\cr
&=&  \Res_{q \to {\bf a}} \Res_{r \to {\bf a}} \Res_{s \to {\bf a}}  {dE_q(p) \over \omega(q)}
{dE_r(q) \over \omega(r)} {dE_s(\qbar) \over \omega(s)} B(r,\overline{r}) B(s, \overline{s})\cr
&&  +\Res_{q \to {\bf a}} \Res_{r \to {\bf a}} \Res_{s \to {\bf a}}  {dE_q(p) \over \omega(q)}
{dE_r(q) \over \omega(r)} {dE_s(\overline{r}) \over \omega(s)} B(r,\qbar) B(s,\overline{s})\cr
&&  +\Res_{q \to {\bf a}} \Res_{r \to {\bf a}} \Res_{s \to {\bf a}}  {dE_q(p) \over \omega(q)}
{dE_r(q) \over \omega(r)} {dE_s(r) \over \omega(s)} \left[ B(\qbar,\overline{r}) B(s,\overline{s})
 \right. \cr
&& \left. + B(\overline{s},\qbar) B(s,\overline{r}) + B(s,\qbar) B(\overline{s},\overline{r}) \right]\cr
&=&
2 \begin{array}{r}
{\epsfxsize 4cm\epsffile{w121.eps}}
\end{array}
 +2 
 \begin{array}{r}
{\epsfxsize 3cm\epsffile{w124.eps}}
\end{array}  
+ \begin{array}{r}
{\epsfxsize 2.8cm\epsffile{w122.eps}}
\end{array}
\eea
where the last expression is obtained using lemma \ref{lemmasymfactor}.


\subsubsection*{Free energy.}

The second order free energy reads
\bea
-2 F^{(2)}
&=&  \Res_{p \to {\bf a}} \Res_{q \to {\bf a}} \Res_{r \to {\bf a}} \Res_{s \to {\bf a}}  {\Phi(p) dE_q(p) \over \omega(q)}
{dE_r(q) \over \omega(r)} {dE_s(\qbar) \over \omega(s)} B(r,\overline{r}) B(s, \overline{s})\cr
&&  +\Res_{p \to {\bf a}} \Res_{q \to {\bf a}} \Res_{r \to {\bf a}} \Res_{s \to {\bf a}}  { \Phi(p) dE_q(p) \over \omega(q)}
{dE_r(q) \over \omega(r)} {\Phi(p) dE_s(\overline{r}) \over \omega(s)} B(r,\qbar) B(s,\overline{s})\cr
&&  +\Res_{p \to {\bf a}} \Res_{q \to {\bf a}} \Res_{r \to {\bf a}} \Res_{s \to {\bf a}}  { \Phi(p) dE_q(p) \over \omega(q)}
{dE_r(q) \over \omega(r)} {dE_s(r) \over \omega(s)} \left[ B(\qbar,\overline{r}) B(s,\overline{s})
 \right. \cr
&& \left. + B(\overline{s},\qbar) B(s,\overline{r}) + B(s,\qbar) B(\overline{s},\overline{r}) \right]\cr
\eea



\subsection{Remark: Teichmuller pant gluings}

Every Riemann surface of genus $g$ with $k$ punctures can be decomposed into 
$2g+k$ pants whose boundaries are $3g+k$ closed geodesics (in the metric with constant negative curvature).
The number of ways (in the combinatorial sense) of gluing $2g+k$ pants by their boundaries is clearly the same as the number of diagrams of ${\cal G}_{k}^{(g)}$, and each diagram corresponds to one pant decomposition.


Example with $k=1$ and $g=2$:
$$
W_1^{(2)}={\epsfxsize 2.2cm\epsffile{teich121.eps}}+2 {\epsfxsize 3cm\epsffile{teich122.eps}}+2 {\epsfxsize 3cm\epsffile{teich123.eps}}
$$




